\section*{Bab \ref{ch:b1:topology} --- Topologi Garis Real}

\begin{solution}{exo:b1:topology:1}
$\intoo{0}{1} \cup \intoo{2}{3}$: \omterm{def:b1:topology:open}{terbuka} (karena gabungan \omterm{def:b1:topology:open}{himpunan terbuka}), dan tak
\omterm{def:b1:topology:closed}{tertutup} (karena $\frac 1n \to 0$ di luarnya).

$\intco{0}{1}$: bukan keduanya. Tak \omterm{def:b1:topology:open}{terbuka} (karena tak ada \omterm{prop:b1:reals:intervals}{selang} di sekitar $0$ di dalamnya);
dan tak \omterm{def:b1:topology:closed}{tertutup} (karena $1 - \frac1n \to 1 \notin$ \omterm{def:b1:logic:sets}{himpunannya}).

$\{0\} \cup \intcc{1}{2}$: \omterm{def:b1:topology:closed}{tertutup} (karena gabungan hingga \omterm{def:b1:topology:closed}{himpunan tertutup}), dan tak
\omterm{def:b1:topology:open}{terbuka} (karena gagal di $0$).

$\R \setminus \Z$: \omterm{def:b1:topology:open}{terbuka} (karena $\Z$ \omterm{def:b1:topology:closed}{tertutup}), dan tak \omterm{def:b1:topology:closed}{tertutup}: sebab barisan
$\bigl(\frac 1n\bigr)$ terletak di dalamnya, tetapi limitnya $0$ termasuk $\Z$,
yakni kabur dari \omterm{def:b1:logic:sets}{himpunannya}.

$\Q \cap \intoo{0}{1}$: bukan keduanya. Tak \omterm{def:b1:topology:open}{terbuka}: karena setiap \omterm{prop:b1:reals:intervals}{selang} di sekitar
bilangan rasional memuat bilangan irasional. Dan tak \omterm{def:b1:topology:closed}{tertutup}: karena ia memuat barisan
yang menuju $\frac{\sqrt 2}{2}$ yang irasional (menurut \omterm{def:b1:topology:dense}{kepadatannya}).
\end{solution}

\begin{solution}{exo:b1:topology:2}
$A = \intoc{0}{1} \cup \{2\}$: di sini $\mathring A = \intoo{0}{1}$,
$\overline A = \intcc{0}{1} \cup \{2\}$, $\partial A = \{0, 1, 2\}$.

$A = \R \setminus \Q$: di sini $\mathring A = \emptyset$ (karena setiap \omterm{prop:b1:reals:intervals}{selang}
memuat bilangan rasional), $\overline A = \R$ (menurut \omterm{def:b1:topology:dense}{kepadatan}
bilangan irasional), $\partial A = \R$.

$A = \bigl\{\frac{(-1)^n n}{n+1}\bigr\}$: suku genapnya menuju $1$,
sedangkan suku ganjilnya ke $-1$, dan keduanya tak termasuk $A$. Jadi $\mathring A =
\emptyset$ (karena titiknya terasing), $\overline A = A \cup \{-1, 1\}$,
$\partial A = \overline A$.
\end{solution}

\begin{solution}{exo:b1:topology:3}
\emph{Lewat komplemennya.} Misalkan $F = \{a_1 < a_2 < \dots < a_k\}$: maka komplemennya
adalah gabungan \omterm{prop:b1:reals:intervals}{selang} \omterm{def:b1:topology:open}{terbuka} $\intoo{-\infty}{a_1}$,
$\intoo{a_i}{a_{i+1}}$, $\intoo{a_k}{+\infty}$ --- yang \omterm{def:b1:topology:open}{terbuka} menurut
\cref{prop:b1:topology:openstable}.

\emph{Lewat barisan.} Misalkan $u_n \in F$, $u_n \to \ell$. Dengan $\varepsilon
= \min\{\abs{a_i - a_j} : i \neq j\}/2 > 0$ (atau sebarang $\varepsilon$ bila
$F$ singleton): di luar suatu peringkat, semua sukunya berada dalam jarak
$\varepsilon$ dari $\ell$, sehingga berjarak $2\varepsilon$ satu sama lain,
yang memaksanya menjadi satu $a_i$ tunggal mulai peringkat itu; lalu
$\ell = a_i \in F$.
\end{solution}

\begin{solution}{exo:b1:topology:4}
$\subseteq$: karena $\overline A \cup \overline B$ \omterm{def:b1:topology:closed}{tertutup} (yaitu gabungan hingga)
dan memuat $A \cup B$, sehingga ia memuat pemuat \omterm{def:b1:topology:closed}{tertutup}
yang \emph{terkecil} $\overline{A \cup B}$. $\supseteq$: karena $A \subseteq A \cup B$
memberikan $\overline A \subseteq \overline{A \cup B}$ (sebab \omterm{def:b1:topology:closure}{penutup} bersifat
monoton: titik lekat $A$ lekat pada \omterm{def:b1:logic:sets}{himpunan} yang lebih besar), dan
demikian pula untuk $B$.

Contoh penyangkal bagi irisannya: $A = \intoo{0}{1}$, $B =
\intoo{1}{2}$: di sini $\overline{A \cap B} = \overline\emptyset = \emptyset$
tetapi $\overline A \cap \overline B = \{1\}$.
\end{solution}

\begin{solution}{exo:b1:topology:5}
Pakailah pencirian lewat barisan
(\cref{thm:b1:topology:seqclosed}). Misalkan $S = \{u_n\} \cup \{\ell\}$
dan $(v_k)$ barisan di $S$ dengan $v_k \to m$; lalu tunjukkan $m \in S$. Ada dua
kasus. Jika suatu nilai $v \in S$ diambil $(v_k)$ tak hingga sering,
maka \omterm{def:b1:seq:subsequence}{subbarisan} yang konstan memberikan $m = v \in S$. Kalau tidak, setiap
nilainya diambil berhingga sering; khususnya, untuk setiap $n$, suku
$u_n$ muncul berhingga sering, dan $\ell$ pun demikian. Maka untuk setiap $N$,
indeks $k$ dengan $v_k \in \{u_0, \dots, u_N, \ell\}$ berhingga
banyaknya: sehingga $v_k$ yang tersisa merupakan suku $u_n$ dengan $n > N$. Diberikan
$\varepsilon > 0$, pilihlah $N$ dengan $\abs{u_n - \ell} \leq
\varepsilon$ untuk $n > N$: maka semua $v_k$ kecuali yang berhingga banyak memenuhi
$\abs{v_k - \ell} \leq \varepsilon$. Jadi $v_k \to \ell$, sehingga $m =
\ell \in S$.
\end{solution}

\begin{solution}{exo:b1:topology:6}
Di sini $U + A = \bigcup_{a \in A} (U + a)$, dan setiap translasinya $U + a$
\omterm{def:b1:topology:open}{terbuka} (karena sertifikat \omterm{prop:b1:reals:intervals}{selangnya} ikut ditranslasikan). Adapun gabungan \omterm{def:b1:topology:open}{himpunan terbuka} bersifat
\omterm{def:b1:topology:open}{terbuka} (\cref{prop:b1:topology:openstable}).

\omterm{def:b1:topology:closed}{Himpunan tertutupnya}: $\Z$ dan $\sqrt 2\,\Z$ \omterm{def:b1:topology:closed}{tertutup} (sebagai $\alpha\Z$:
karena barisan yang konvergen akhirnya konstan, bandingkan
\cref{ex:b1:topology:closed}). Adapun jumlahnya $\Z + \sqrt 2\,\Z$ bersifat \omterm{def:b1:topology:dense}{padat}
di $\R$ (\cref{exo:b1:reals:9}) tetapi bukan $\R$ (karena ia terbilang, atau
sederhananya: $\frac{\sqrt 2}{2} \notin \Z + \sqrt2\Z$, sebab kalau tidak $\sqrt 2$
akan rasional --- karena menulis $\frac{\sqrt2}{2} = m + n\sqrt 2$
memaksa $(2n - 1)\sqrt 2 = -2m$, sehingga $\sqrt 2 \in \Q$ kecuali bila $n =
\frac12$, yang mustahil). Adapun \omterm{def:b1:logic:sets}{himpunan} bagian sejati yang \omterm{def:b1:topology:dense}{padat} tak \omterm{def:b1:topology:closed}{tertutup}: karena
\omterm{def:b1:topology:closure}{penutupnya} adalah $\R \neq$ dirinya sendiri.
\end{solution}

\begin{solution}{exo:b1:topology:7}
$\Z$: \omterm{def:b1:topology:open}{persekitaran} $\intoo{n - \frac12}{n + \frac12}$ bagi $n$
berpotongan dengan $\Z$ hanya pada $n$.

$A = \{\frac1n : n \in \N^*\}$: di sekitar $\frac 1n$, \omterm{prop:b1:reals:intervals}{selang} berjari-jari
$\frac{1}{n} - \frac{1}{n+1} = \frac{1}{n(n+1)}$ (yang diparuhkan, misalnya)
mengasingkannya dari tetangganya --- jadi semua titiknya terasing. Namun $0 \in
\overline A \setminus A$: sehingga titik yang terasing tak mencegah orang luar
yang lekat.

Jika semua titik $A$ terasing: maka tak ada titik $A$ yang \omterm{def:b1:topology:closure}{interior}, karena
sebuah titik \omterm{def:b1:topology:closure}{interior} mempunyai seluruh \omterm{prop:b1:reals:intervals}{selang} berisi tetangga $A$ di sekitarnya
(dan \omterm{prop:b1:reals:intervals}{selang} bersifat tak hingga), yang bertentangan dengan keterasingannya. Jadi $\mathring A
= \emptyset$.
\end{solution}

\begin{solution}{exo:b1:topology:8}
Jika $A \subseteq \intcc{-M}{M}$, maka \omterm{def:b1:topology:closed}{himpunan tertutup} $\intcc{-M}{M}$
memuat $A$, sehingga ia memuat $\overline A$ (yaitu pemuat \omterm{def:b1:topology:closed}{tertutup}
yang terkecil): jadi $\overline A$ terbatas.

Misalkan $s = \sup A$ (yang hingga). Karena $A \subseteq \overline A$, maka $\sup
\overline A \geq s$. Sebaliknya, $\overline A \subseteq
\intoc{-\infty}{s}$: sebab setengah garisnya \omterm{def:b1:topology:closed}{tertutup} dan memuat $A$; sehingga
setiap unsur $\overline A$ bernilai $\leq s$, yang memberikan $\sup\overline A
\leq s$. Jadi sama.

Adapun $s \in \overline A$: menurut pencirian lewat $\varepsilon$
(\cref{prop:b1:reals:epsilon}), setiap \omterm{prop:b1:reals:intervals}{selang} $\intoo{s -
\varepsilon}{s + \varepsilon}$ memuat sebuah unsur $A$: jadi $s$
bersifat lekat.
\end{solution}

\begin{solution}{exo:b1:topology:9}
Andaikan $A$ \omterm{def:b1:topology:open}{terbuka} dan \omterm{def:b1:topology:closed}{tertutup}, dengan $a \in A$ dan $b \in \R
\setminus A$; tanpa mengurangi keumumannya $a < b$. \omterm{def:b1:logic:sets}{Himpunan} $B = A \cap
\intcc{a}{b}$ tak kosong (karena $a$), dan terbatas: misalkan $s = \sup B$. Menurut
\cref{exo:b1:topology:8}, $s \in \overline B \subseteq \overline A =
A$ (karena $A$ \omterm{def:b1:topology:closed}{tertutup}). Perhatikan $s \leq b$, dan karena $b \notin A$: maka $s < b$.
Sekarang $A$ \omterm{def:b1:topology:open}{terbuka}: sehingga suatu \omterm{prop:b1:reals:intervals}{selang} $\intoo{s - r}{s + r}$ terletak di $A$,
dan kita boleh mengambil $r < b - s$. Maka $s + \frac{r}{2}$ termasuk $A
\cap \intcc{a}{b} = B$ dan melampaui $s$ --- yang bertentangan dengan $s = \sup
B$. Jadi tak ada pasangan $(a, b)$ yang demikian: sehingga salah satu dari $A$ dan $\R \setminus A$
bersifat kosong.
\end{solution}

\begin{solution}{exo:b1:topology:10}
\omterm{def:b1:logic:sets}{Himpunan} $I_x$ merupakan gabungan \omterm{prop:b1:reals:intervals}{selang} \omterm{def:b1:topology:open}{terbuka} yang semuanya memuat $x$: sehingga ia \omterm{def:b1:topology:open}{terbuka},
dan ia sebuah \omterm{prop:b1:reals:intervals}{selang}, karena cembung --- sebab jika $u < z < v$ dengan $u, v \in
I_x$, maka $u$ dan $v$ terletak pada subselang \omterm{def:b1:topology:open}{terbuka} $J_u \ni x$, $J_v \ni
x$ dari $U$, dan $J_u \cup J_v$ merupakan \omterm{prop:b1:reals:intervals}{selang} (karena keduanya memuat $x$)
di dalam $U$ yang memuat $z$; sehingga $z \in I_x$
(\cref{prop:b1:reals:intervals}).

Jika $I_x \cap I_y \neq \emptyset$: maka $I_x \cup I_y$ merupakan \omterm{prop:b1:reals:intervals}{selang}
\omterm{def:b1:topology:open}{terbuka} (yang cembung, karena dua \omterm{prop:b1:reals:intervals}{selang} yang bertindihan) yang termuat $U$
dan memuat $x$ maupun $y$, sehingga $I_x \cup I_y \subseteq I_x$ dan
$\subseteq I_y$ menurut kemaksimalan masing-masingnya: jadi $I_x = I_y$.

Jadi $U$ merupakan gabungan lepas atas \omterm{def:b1:logic:sets}{himpunan} $I_x$ yang berbeda-beda (karena setiap $x \in
U$ terletak di $I_x$-nya sendiri). Keterbilangannya: setiap \omterm{prop:b1:reals:intervals}{selang} \omterm{def:b1:topology:open}{terbuka}
$I$ yang tak kosong pada keluarganya memuat sebuah bilangan rasional $q_I$
(\cref{thm:b1:reals:density}), dan \omterm{prop:b1:reals:intervals}{selang} lepas yang berbeda memperoleh
bilangan rasional yang berbeda: sehingga keluarganya terinjeksi ke $\Q$, yang terbilang
(karena ia terindeks oleh pasangan bilangan bulat). Jadi \omterm{prop:b1:reals:intervals}{selangnya} paling banyak
terbilang.
\end{solution}

\begin{solution}{exo:b1:topology:11}
Klaimnya $\overline A = A \cup A'$. ($\supseteq$) Selalu $A \subseteq \overline
A$; dan jika $x \in A'$, maka setiap \omterm{def:b1:topology:open}{persekitaran} $x$ berpotongan dengan $A
\setminus \{x\} \subseteq A$, sehingga $x$ bersifat lekat. ($\subseteq$)
Misalkan $x \in \overline A$. Jika $x \in A$, selesai. Jika $x \notin A$, maka
setiap \omterm{def:b1:topology:open}{persekitaran} $x$ berpotongan dengan $A = A \setminus \{x\}$: jadi $x \in
A'$.

Akibatnya $A$ \omterm{def:b1:topology:closed}{tertutup} $\iff$ $A = \overline A = A \cup A'$
$\iff$ $A' \subseteq A$.

$A = \{\frac 1n\}$: di sini $0$ merupakan titik akumulasi (karena $\frac 1n \to
0$, dengan suku $\neq 0$); sedangkan setiap $\frac 1n$ terasing
(\cref{exo:b1:topology:7}), sehingga tak termasuk $A'$; dan sebuah titik $x
\notin A \cup \{0\}$ mempunyai seluruh \omterm{prop:b1:reals:intervals}{selang} yang menghindari $A$ (yaitu di antara
kedua tetangga $x$ pada $A \cup \{0\}$, atau di luar $1$).
Jadi $A' = \{0\}$.

$\Z' = \emptyset$: karena setiap bilangan bulat terasing, dan setiap yang bukan bilangan bulat
mempunyai \omterm{def:b1:topology:open}{persekitaran} di dalam $\R \setminus \Z$.

$\Q' = \R$: karena setiap \omterm{prop:b1:reals:intervals}{selang} di sekitar sebarang bilangan real memuat tak hingga
banyak bilangan rasional (\cref{thm:b1:reals:density}), khususnya satu
yang berbeda dari pusatnya.
\end{solution}

\begin{solution}{exo:b1:topology:12}
\begin{enumerate}
  \item Untuk setiap $a \in A$: $\abs{x - a} \leq \abs{x - y} +
        \abs{y - a}$, sehingga $d_A(x) \leq \abs{x - y} + \abs{y - a}$;
        lalu mengambil \omterm{def:b1:reals:bounds}{infimum} atas $a$: $d_A(x) \leq \abs{x - y} +
        d_A(y)$. Menukar $x$ dan $y$ memberikan ketaksamaan
        yang lain: $\abs{d_A(x) - d_A(y)} \leq \abs{x - y}$.
  \item $d_A(x) = 0$ $\iff$ untuk setiap $\varepsilon > 0$ ada
        $a \in A$ dengan $\abs{x - a} < \varepsilon$ $\iff$ setiap
        \omterm{prop:b1:reals:intervals}{selang} di sekitar $x$ berpotongan dengan $A$ $\iff$ $x \in \overline A$.
        Jika $F$ \omterm{def:b1:topology:closed}{tertutup} dan $x \notin F = \overline F$, maka
        $d_F(x) \neq 0$, yakni $d_F(x) > 0$.
  \item Jika $d_F(x) < \varepsilon$, tetapkan $r = \varepsilon -
        d_F(x) > 0$: maka untuk $\abs{y - x} < r$, bagian (1) memberikan
        $d_F(y) \leq d_F(x) + \abs{x - y} < \varepsilon$: jadi
        \omterm{prop:b1:reals:intervals}{selang} $\intoo{x - r}{x + r}$ terletak di $V_\varepsilon$,
        sehingga ia \omterm{def:b1:topology:open}{terbuka}; dan ia memuat $F$ karena $d_F = 0$
        di sana. Akhirnya $x \in \bigcap_{\varepsilon>0}
        V_\varepsilon$ $\iff$ $d_F(x) < \varepsilon$ untuk setiap
        $\varepsilon$ $\iff$ $d_F(x) = 0$ $\iff$ $x \in \overline
        F = F$. Dan karena $\bigcap_{\varepsilon > 0} V_\varepsilon =
        \bigcap_{n \geq 1} V_{1/n}$, setiap \omterm{def:b1:topology:closed}{himpunan tertutup} merupakan
        irisan terbilang atas \omterm{def:b1:topology:open}{himpunan terbuka}.
\end{enumerate}
\end{solution}

\begin{solution}{pb:b1:topology:1}
\textbf{1.} Di sini $C_1 = \intcc{0}{\frac13} \cup \intcc{\frac23}{1}$
dan
\[
C_2 = \intcc{0}{\tfrac19} \cup \intcc{\tfrac29}{\tfrac13}
\cup \intcc{\tfrac23}{\tfrac79} \cup \intcc{\tfrac89}{1} .
\]
Lewat induksi: jika $C_n$ merupakan gabungan lepas atas $2^n$ ruas \omterm{def:b1:topology:closed}{tertutup}
yang berpanjang $3^{-n}$, maka menghapus sepertiga tengahnya yang \omterm{def:b1:topology:open}{terbuka} dari masing-masingnya
menyisakan dua ruas \omterm{def:b1:topology:closed}{tertutup} berpanjang $3^{-n-1}$ per induknya:
sehingga $2^{n+1}$ ruas, yang saling lepas (karena anak dari induk yang berbeda
terpisah sebab induknya sudah terpisah; sedangkan anak dari satu
induk terpisah oleh celah yang dibuang).

\textbf{2.} Setiap $C_n$ merupakan gabungan hingga atas ruas, sehingga
\omterm{def:b1:topology:closed}{tertutup}; lalu $C = \bigcap C_n$ merupakan irisan \omterm{def:b1:topology:closed}{himpunan tertutup}:
jadi \omterm{def:b1:topology:closed}{tertutup} (\cref{def:b1:topology:closed}); dan terbatas (karena $\subseteq
\intcc{0}{1}$): sehingga kompak menurut \cref{thm:b1:topology:compact}.
Tak kosong: karena $0$ terletak di ruas paling kiri setiap $C_n$. Misalkan
$a$ titik ujung sebuah ruas $S$ pada $C_n$. Untuk $m \leq n$,
$a \in C_n \subseteq C_m$. Adapun untuk tahap berikutnya: penghapusan sepertiga
tengahnya tak pernah membuang titik ujung, dan $a$ kembali menjadi titik ujung
salah satu dari kedua anak $S$ (yaitu anak yang menyentuh $a$); sehingga lewat
induksi $a \in C_m$ untuk setiap $m \geq n$: jadi $a \in C$.

\textbf{3.} Panjang total $C_n$: $2^n \cdot 3^{-n} =
(\frac23)^n \to 0$. Diberikan $\varepsilon > 0$, pilihlah $n$ dengan
$(\frac23)^n \leq \varepsilon$: maka $C \subseteq C_n$, yaitu gabungan
atas ruas yang berhingga banyak dengan panjang total $\leq \varepsilon$.

\textbf{4.} Misalkan $I \subseteq C$ sebuah \omterm{prop:b1:reals:intervals}{selang} dengan dua titik
yang berbeda. Untuk setiap $n$: $I \subseteq C_n$, dan $I$, karena cembung,
harus terletak di dalam \emph{satu} ruas $C_n$ (karena berpotongan dengan dua
ruas akan memaksa $I$ memuat sebuah titik celah di antara
keduanya, yang berada di luar $C_n$). Sehingga panjang $I$ bernilai $\leq
3^{-n}$ untuk setiap $n$: yang bertentangan. Jadi satu-satunya \omterm{prop:b1:reals:intervals}{selang}
di dalam $C$ adalah yang kosong atau singleton; khususnya tak ada
$\intoo{x-r}{x+r}$ yang muat di dalam $C$: jadi $\mathring C = \emptyset$.
Dan karena $C$ \omterm{def:b1:topology:closed}{tertutup}, $\overline C = C$ berinterior kosong: sehingga $C$
tak \omterm{def:b1:topology:dense}{padat} di mana pun.

\textbf{5.} Tulislah $\varphi_0(x) = \frac x3$ dan $\varphi_2(x) =
\frac{2 + x}{3}$, yaitu bijeksi afin yang naik dari $\intcc{0}{1}$
pada $\intcc{0}{\frac13}$ dan $\intcc{\frac23}{1}$. Klaimnya:
$C_{n+1} = \varphi_0(C_n) \cup \varphi_2(C_n)$. Untuk $n = 0$ inilah
pertanyaan 1. Lewat induksi: \omterm{def:b1:logic:map}{pemetaan} afin yang naik mengirim
sepertiga tengah sebuah ruas ke sepertiga tengah ruas
petanya, sehingga menghapus sepertiga tengah berkomutasi dengan $\varphi_0$ dan
$\varphi_2$; lalu menerapkan langkah penghapusannya pada $C_{n+1} =
\varphi_0(C_n) \cup \varphi_2(C_n)$ menghasilkan $C_{n+2} =
\varphi_0(C_{n+1}) \cup \varphi_2(C_{n+1})$. Lalu mengiriskan atas
$n$: untuk $x \leq \frac13$, $x \in C \iff x \in \varphi_0(C_n)$
untuk setiap $n$ $\iff 3x \in \bigcap C_n = C$; demikian pula pada
$\intcc{\frac23}{1}$; dan tak ada titik $\intoo{\frac13}{\frac23}$ yang
terletak di $C_1$. Jadi $C = \varphi_0(C) \cup \varphi_2(C)$,
secara lepas.

\textbf{6.} Lewat induksi pada $n$; adapun kasus $n = 0$ mengatakan bahwa setiap $x \in
\intcc{0}{1}$ mempunyai sandi, dan itulah \cref{pb:b1:reals:1}
(pertanyaan 9 untuk $x < 1$; sedangkan $1 = (0.\overline 2)_3$). Andaikan
kesetaraannya berlaku pada peringkat $n$. Jika $x \in C_{n+1}$: maka menurut pertanyaan 5, $x =
\varphi_i(z)$ dengan $z \in C_n$ dan $i \in \{0, 2\}$; lalu jika $(e_k)$
sandi $z$ yang $n$ angka pertamanya bebas-$1$, maka $(i, e_1,
e_2, \dots)$ berjumlah parsial $\frac i3 + \frac13\sum_{k\leq m}
e_k 3^{-k} \to \varphi_i(z) = x$: yaitu sandi $x$ dengan $n +
1$ angka pertamanya bebas-$1$. Sebaliknya, jika $x$ mempunyai sandi $(d_k)$ dengan
$d_1, \dots, d_{n+1} \in \{0, 2\}$: maka untaian tergesernya $(d_2,
d_3, \dots)$ bernilai suatu $z \in \intcc{0}{1}$, yang $n$ angka
pertamanya bebas-$1$, dan perhitungan jumlah parsialnya yang dibaca
mundur memberikan $x = \varphi_{d_1}(z)$; lalu lewat induksi $z \in
C_n$, sehingga $x \in C_{n+1}$ menurut pertanyaan 5. Akhirnya: sandi
yang seluruhnya bebas-$1$ menaruh $x$ di setiap $C_n$, sehingga di $C$; sebaliknya
jika $x \in C$, maka untuk setiap $n$ salah satu dari paling banyak dua sandi
$x$ (\cref{pb:b1:reals:1}, pertanyaan 11) mempunyai $n$ angka pertamanya
bebas-$1$; lalu satu sandi yang tetap pasti berlaku untuk $n$ yang sebesar-besarnya
(menurut sarang merpati di antara kedua sandinya), dan sandi yang $n$ angka
pertamanya bebas-$1$ untuk $n$ yang sebesar-besarnya bersifat bebas-$1$
sepenuhnya.

\textbf{7.} Di sini $0 = (0.\overline 0)_3$, $1 = (0.\overline 2)_3$,
$\frac13 = (0.0\overline{2})_3$ (yaitu kembaran tak sejati dari $(0.1)_3$),
dan $\frac23 = (0.2\overline{0})_3$. Lewat jumlah geometrinya:
\[
(0.\overline{02})_3 = \sum_{j\geq1} \frac{2}{9^{\,j}}
= \frac{2/9}{1 - 1/9} = \frac14 ,
\qquad
(0.\overline{20})_3 = \sum_{j\geq1} \frac{2}{3\cdot 9^{\,j-1}}
= \frac{2/3}{1 - 1/9} = \frac34 ,
\]
keduanya bebas-$1$: jadi $\frac14, \frac34 \in C$. (Adapun jumlah tak hingganya
menyingkat \omterm{def:b1:reals:bounds}{supremum} jumlah parsialnya, seperti pada \cref{pb:b1:reals:1}.)
Titik ujung ruas $C_n$ berbentuk $m/3^n$
(lewat induksi: karena titik ujung anaknya adalah titik ujung induknya atau berbeda
dari salah satunya sejauh kelipatan $3^{-n-1}$). Jika $\frac14 =
\frac{m}{3^n}$ maka $3^n = 4m$, padahal $4 \nmid 3^n$: yang mustahil.
Jadi $\frac14 \in C$ tanpa pernah menjadi titik ujung.

\textbf{8.} Andaikan $x$ mempunyai dua sandi bebas-$1$ yang berbeda. Mempunyai
dua sandi sama sekali berarti (menurut \cref{pb:b1:reals:1}, pertanyaan 11, basis
$3$) bahwa $x = m/3^N \in \intoo{0}{1}$ dan kedua sandinya adalah:
yang berakhir, dengan angka taknol terakhirnya $d_N \in \{1, 2\}$
lalu disusul $0$, beserta kembarannya, dengan $d_N - 1$ pada posisi $N$
lalu disusul $2$. Jika $d_N = 1$ maka yang pertama memuat $1$; sedangkan jika $d_N
= 2$ maka kembarannya membawa $d_N - 1 = 1$. Bagaimanapun paling banyak satu
dari pasangannya bebas-$1$: yang bertentangan. Jadi setiap $x \in C$ mempunyai
tepat satu sandi bebas-$1$ (dengan keberadaannya menurut pertanyaan 6), dan
untaian $\{0,2\}$ yang berbeda bernilai berbeda. Adapun setiap
untaian $\{0,2\}$ bernilai di $\intcc{0}{1}$ (karena jumlah parsialnya $\leq
1$) dengan semua awalannya bebas-$1$, sehingga nilainya berada di setiap $C_n$,
yakni di $C$: jadi \omterm{def:b1:logic:map}{pemetaan} nilainya merupakan bijeksi dari
untaian $\{0,2\}$ pada $C$.

\textbf{9.} Misalkan $(d^{(k)})$ sandi bebas-$1$ bagi $x_k$ lalu
tetapkan $e_k = 2 - d^{(k)}_k \in \{0, 2\}$: yaitu untaian $\{0,2\}$ yang
nilainya $y$ terletak di $C$ dan mempunyai $(e_k)$ sebagai sandi bebas-$1$
tunggalnya (pertanyaan 8). Untuk setiap $k$ sandi $y$ dan $x_k$
berbeda pada posisi $k$, sehingga $y \neq x_k$: jadi tak ada \omterm{def:b1:logic:map}{pemetaan} $\N^* \to C$ yang
\omterm{def:b1:logic:inj}{surjektif}. Yaitu \omterm{def:b1:logic:sets}{himpunan} takterbilang yang berpanjang nol: besar menurut
kardinalitasnya, kecil menurut ukurannya --- serentak.

\textbf{10.} Kekompakannya: yaitu ketertutupan irisan tak hingganya
ditambah keterbatasannya (pertanyaan 2) --- itulah satu sifat
yang tak dibawa oleh satu pemuatan tunggal. Panjang nolnya: dari $C \subseteq
C_n$ yang berpanjang total $(\frac23)^n$ (pertanyaan 3). Ketakpadatannya
di mana pun: karena $C \subseteq C_n$ memaksa \omterm{prop:b1:reals:intervals}{selang} di dalam $C$ berpanjang
$\leq 3^{-n}$ (pertanyaan 4). Sedangkan ketakterbilangannya sama sekali tak menumpang pada
pemuatan: ia memerlukan seluruh struktur irisannya, yang
tersandikan pada bijeksi pertanyaan 8.

\textbf{11.} Balikkan $d_n$ menjadi $2 - d_n$: maka untaian barunya tetap
untaian $\{0,2\}$, sehingga nilainya $x_n$ terletak di $C$; dan jumlah
parsial di luar peringkat $n$ berselisih persis $2\cdot3^{-n}$, sehingga
$\abs{x_n - x} = 2\cdot3^{-n}$. Jadi $x_n \neq x$ dan $x_n \to
x$: sehingga setiap titik $C$ merupakan limit titik lain pada $C$ --- jadi $C$
bersifat \emph{sempurna}, tanpa titik yang terasing.

\textbf{12.} Menurut induksi pertanyaan 5, ruas
$C_n$ tepat berupa $\intcc{t}{t + 3^{-n}}$ dengan $t$ menjelajahi
nilai untaian berpanjang $n$ atas $\{0,2\}$. Diberikan $x \in C$
bersandi $(d_k)$, maka pemenggalannya $t_n$ (yaitu angka $d_1 \dots d_n$
lalu $0$) merupakan titik ujung kiri, dan $0 \leq x - t_n
\leq 3^{-n}$: jadi titik ujungnya \omterm{def:b1:topology:dense}{padat} di $C$. Mereka membentuk \omterm{def:b1:logic:sets}{himpunan} bagian
$\{m/3^n : m, n\}$, yaitu \omterm{def:b1:logic:sets}{himpunan} yang terindeks oleh pasangan bilangan bulat, sehingga
terbilang (seperti untuk $\Q$ pada \cref{exo:b1:topology:10}). Dan karena $C$
takterbilang (pertanyaan 9), semua titik $C$ kecuali yang terbilang banyaknya
bukan titik ujung --- adapun $\frac14$ (pertanyaan 7) adalah pucuk
gunung es itu yang kasatmata.

\textbf{13.} Pilihlah $n$ dengan $3^{-n} < y - x$. Keduanya $x, y \in
C_n$, dan keduanya tak dapat terletak pada ruas yang sama (karena panjangnya $3^{-n} <
y - x$): sehingga celah yang dibuang di antara ruasnya menyediakan $z$
dengan $x < z < y$ dan $z \notin C_n \supseteq C$. Jadi sebarang dua
titik $C$ terpisah oleh komplemennya: sehingga satu-satunya \omterm{def:b1:logic:sets}{himpunan} bagian
$C$ yang cembung adalah singleton --- jadi $C$ terputus seluruhnya.

\textbf{14.} Jika $(d_k)$ sandi bebas-$1$ bagi $x$, maka untaian
$(2 - d_k)$ kembali menjadi untaian $\{0,2\}$, dengan jumlah parsial
\[
\sum_{k=1}^{n} (2 - d_k)3^{-k} = (1 - 3^{-n}) - \sum_{k=1}^n
d_k 3^{-k} \longrightarrow 1 - x :
\]
sehingga $1 - x \in C$. Jadi $1 - C \subseteq C$, lalu menerapkan petanya
dua kali memberikan $1 - C = C$: sehingga \omterm{pb:b1:topology:1}{himpunan Cantor} setangkup terhadap
$\frac12$.

\textbf{15.} Jumlah parsialnya $t_n \to x$ dan $t'_n \to x'$
(karena barisan yang naik konvergen ke \omterm{def:b1:reals:bounds}{supremumnya}, yakni
nilainya), sehingga $t_n + t'_n \to x + x'$ menurut
\cref{thm:b1:seq:operations}. Diberikan $y \in \intcc{0}{1}$ dengan
sandi $(e_k)$, pilihlah $(a_k, b_k) = (0,0), (0,2), (2,2)$
menurut $e_k = 0, 1, 2$: maka $a_k + b_k = 2e_k$, sehingga
untaian $(a_k)$, $(b_k)$ merupakan untaian $\{0,2\}$ dengan nilai $x,
x' \in C$, dan
\[
x + x' = \lim_n\,(t_n + t'_n) = \lim_n 2\sum_{k=1}^n e_k 3^{-k}
= 2y :
\]
jadi setiap $y \in \intcc{0}{1}$ merupakan titik tengah dua titik $C$.

\textbf{16.} Pertanyaan 15 memberikan $\intcc{0}{2} = 2\,\intcc{0}{1}
\subseteq C + C$, dan $C + C \subseteq \intcc{0}{1} +
\intcc{0}{1} = \intcc{0}{2}$: jadi sama. Lalu, dengan memakai $1 - C = C$:
\[
C - C = C + (C - 1) = (C + C) - 1 = \intcc{-1}{1} .
\]
Contoh konkretnya: $1 = \frac14 + \frac34$, yaitu jumlah dua
anggota $C$ yang bukan titik ujung.

\textbf{17.} Panjang mengukur seberapa banyak garisnya yang ditempati \omterm{def:b1:logic:sets}{himpunannya}
sendiri; ia tak mengatakan apa pun tentang \omterm{def:b1:logic:sets}{himpunan} \emph{jumlahnya}, yang merupakan
peta dari keluarga berparameter dua $C \times C$ di bawah $(x,
x') \mapsto x + x'$ --- karena kedua untaian angkanya dipilih
secara saling bebas, dan kebebasan itulah yang persis mengisi
$\intcc{0}{2}$. Tak ada teorema yang membatasi panjang \omterm{def:b1:logic:sets}{himpunan} jumlah oleh
panjang penjumlahnya, dan $C$ adalah buktinya bahwa tak akan pernah ada.

\textbf{18.} Menurut \cref{pb:b1:reals:1} (pertanyaan 18), $x$
rasional jika dan hanya jika ekspansi sejatinya periodik pada akhirnya. Adapun
sandi bebas-$1$ bagi $x \in C$ entah ekspansi sejati itu entah
kembaran tak sejati dari yang berakhir; dan untaian yang berakhir beserta
kembarannya (yang akhirnya $2$ terus) sama-sama periodik pada
akhirnya, sehingga keperiodikan sandi bebas-$1$-nya setara dengan
kerasionalan $x$. Pembagian panjang $\frac1{13}$ dalam basis $3$
(dengan $r_0 = 1$): $3 = 13\cdot0 + 3$, $9 = 13\cdot0 + 9$, $27 =
13\cdot2 + 1$, lalu sisanya kembali ke $1$: sehingga angkanya
$\overline{002}$, jadi $\frac1{13} = (0.\overline{002})_3$,
yang bebas-$1$ dan periodik: yaitu anggota $C$ yang rasional. (Periksa:
$\frac{2/27}{1 - 1/27} = \frac{2}{26} = \frac1{13}$.)

\textbf{19.} Untaian dengan $d_k = 2$ pada posisi
segitiga $k = \frac{j(j+1)}{2}$ dan $0$ di tempat lain merupakan
untaian $\{0,2\}$, sehingga nilainya $x^*$ termasuk $C$ (pertanyaan
8). Ia mempunyai tak hingga banyak $2$ dengan jarak $j + 1 \to \infty$
di antara yang berurutan, sehingga ia tak periodik pada akhirnya (karena
periode $T$ akhirnya akan memaksa $2$ berjarak $\leq T$: yaitu argumen
jarak yang tumbuh pada \cref{pb:b1:reals:1}, pertanyaan 20); lalu menurut
pertanyaan 18, $x^* \notin \Q$. Dan menurut pertanyaan 9 ditambah
keterbilangan $\Q$, semua anggota $C$ kecuali yang terbilang banyaknya bersifat
irasional: jadi $x^*$ itu lazim, bukan kekecualian.

\textbf{20.} Misalkan $y \in \intcc{0}{1}$: ia mempunyai sandi biner
$(c_k)$ dengan $c_k \in \{0, 1\}$ (\cref{pb:b1:reals:1}, pertanyaan
9, basis $2$; adapun $y = 1$ memakai untaian yang seluruhnya $1$). Maka $(2c_k)$
merupakan untaian $\{0,2\}$, nilainya $x$ terletak di $C$, dan $h(x)$ adalah
nilai $(c_k)$, yakni $y$: jadi $h$ memetakan $C$ pada
$\intcc{0}{1}$. Seandainya $C$ merupakan peta dari sebuah \omterm{def:b1:logic:map}{pemetaan} dari $\N^*$,
maka menyusunnya dengan $h$ akan mendaftar seluruh $\intco{0}{1}$,
yang bertentangan dengan teorema diagonal pada \cref{pb:b1:reals:1}
(pertanyaan 22): jadi $C$ takterbilang, sekali lagi. Yaitu \omterm{def:b1:logic:sets}{himpunan} berpanjang nol
yang \omterm{def:b1:logic:inj}{surjektif} pada sebuah ruas yang penuh.

\textbf{21.} Menurut pertanyaan 5, $C$ merupakan gabungan lepas atas
$\varphi_0(C)$ dan $\varphi_2(C)$, yang masing-masingnya translasi dari
salinan terskala $\frac13 C$. Lalu keaditifan, penskalaan dan kekekalan terhadap
translasi memberikan
\[
L(C) = L(\varphi_0(C)) + L(\varphi_2(C))
= \tfrac13 L(C) + \tfrac13 L(C) = \tfrac23\,L(C),
\]
sehingga $\frac13 L(C) = 0$: jadi $L(C) = 0$. Keserupaan diri saja sudah
menghukum $C$ menjadi berpanjang nol --- dan pertanyaan 3 hanya melaksanakan
hukumannya.

\textbf{22.} Sebuah ruas berpanjang $l_n$ kehilangan \omterm{prop:b1:reals:intervals}{selang} di tengahnya
yang berpanjang $4^{-(n+1)}$, sehingga menyisakan dua ruas berpanjang $l_{n+1}
= \frac{l_n - 4^{-(n+1)}}{2}$; lalu dari $l_0 = 1$, induksi
menegaskan $l_n = \frac{2^n + 1}{2\cdot4^n}$: memang
$\frac12\Bigl(\frac{2^n+1}{2\cdot4^n} - \frac{1}{4^{n+1}}\Bigr)
= \frac{2(2^n + 1) - 1}{2\cdot4^{n+1}} = \frac{2^{n+1} +
1}{2\cdot4^{n+1}}$, dan $l_n > 0$ selalu: jadi konstruksinya tak pernah
kelaparan. Adapun $K = \bigcap K_n$ \omterm{def:b1:topology:closed}{tertutup} dan terbatas, sehingga kompak;
dan \omterm{prop:b1:reals:intervals}{selang} di dalam $K$ terletak pada satu ruas $K_n$, yang berpanjang
$l_n \to 0$: jadi \omterm{def:b1:topology:closure}{interiornya} kosong. Panjang yang terbuang: $\sum_{n\geq0} 2^n
\cdot 4^{-(n+1)} = \frac14\sum_{n\geq0}\bigl(\frac12\bigr)^n =
\frac12$, dan setiap $K_n$ berpanjang total $2^n l_n = \frac{2^n +
1}{2^{n+1}} > \frac12$. Sekarang misalkan \omterm{prop:b1:reals:intervals}{selang} \omterm{def:b1:topology:open}{terbuka} yang berhingga banyak
bergabungan $U \supseteq K$. Menurut \omterm{def:b1:logic:statement}{pernyataan} pendamping pada
\cref{ex:b1:topology:nested}, $U \supseteq K_n$ untuk suatu $n$;
lalu dengan menerima keaditifan panjang atas gabungan hingga \omterm{prop:b1:reals:intervals}{selang},
panjang total \omterm{prop:b1:reals:intervals}{selang} penyelimutnya sekurang-kurangnya sebesar panjang
$K_n$, yang melampaui $\frac12$. Jadi $K$ tak \omterm{def:b1:topology:dense}{padat} di mana pun, namun tak ada
selimut yang murah: sehingga kekecilan topologi (tak \omterm{def:b1:topology:dense}{padat} di mana pun) dan
kekecilan metrik (berpanjang nol) sungguh gagasan yang berbeda,
dan $K$ memisahkan keduanya.

\textbf{23.} Ketercapaiannya: misalkan $d = d(x, F)$ lalu pilihlah $a_k \in F$
dengan $\abs{x - a_k} \leq d + \frac1k$: maka $a_k$ terbatas, sehingga
Bolzano--Weierstrass (\cref{thm:b1:seq:bw}) mengekstrak
$a_{\varphi(k)} \to a$, dengan $a \in F$ (karena $F$ \omterm{def:b1:topology:closed}{tertutup},
\cref{thm:b1:topology:seqclosed}) dan $\abs{x - a} = \lim
\abs{x - a_{\varphi(k)}} = d$. Sekarang maksimumnya: jika $y \in C$, maka
$d(y, C) = 0$; kalau tidak, $y$ terletak pada celah yang dibuang pada suatu tahap
$n \geq 1$, yaitu \omterm{prop:b1:reals:intervals}{selang} \omterm{def:b1:topology:open}{terbuka} berpanjang $3^{-n}$ yang kedua
titik ujungnya termasuk $C$ (pertanyaan 2), sehingga $d(y, C) \leq
\frac{3^{-n}}{2} \leq \frac16$, dengan kesamaannya menuntut $n = 1$
dan $y$ pada pusat celah $\intoo{\frac13}{\frac23}$,
yakni $y = \frac12$; dan memang $d\bigl(\frac12, C\bigr) =
\frac16$ karena $C \cap \intoo{\frac13}{\frac23} = \emptyset$ dan
$\frac13, \frac23 \in C$. Jadi $\max_{y\in\intcc{0}{1}} d(y, C)
= \frac16$, yang tercapai persis di $\frac12$.

\textbf{24.} Titik ujungnya membentuk \omterm{def:b1:logic:sets}{himpunan} bagian $C$ yang terbilang dan \omterm{def:b1:topology:dense}{padat}
(pertanyaan 12): daftarkanlah mereka sebagai satu barisan tunggal $(e_j)_{j\geq1}$,
yaitu barisan di $C$. Limit \omterm{def:b1:seq:subsequence}{subbarisannya} semuanya terletak di $C$ (karena $C$
\omterm{def:b1:topology:closed}{tertutup}). Sebaliknya, tetapkan $x \in C$: untuk setiap $n$, ruas
$C_m$ yang memuat $x$ (dengan $m \geq n$) mempunyai titik ujungnya
dalam jarak $3^{-m} \leq 3^{-n}$ dari $x$, sehingga tak hingga banyak titik ujung
yang berbeda terletak dalam jarak $3^{-n}$ dari $x$; lalu pilihlah indeks $j_1 < j_2
< \dots$ dengan $\abs{e_{j_n} - x} \leq 3^{-n}$: yaitu \omterm{def:b1:seq:subsequence}{subbarisan}
yang konvergen ke $x$. Jadi \omterm{def:b1:logic:sets}{himpunan} limit \omterm{def:b1:seq:subsequence}{subbarisan} $(e_j)$
tepat sama dengan $C$ --- yaitu satu barisan yang menggerombol pada titik yang
takterbilang banyaknya, yakni ujung yang berlawanan dari barisan yang konvergen, yang
\omterm{def:b1:logic:sets}{himpunan} gerombolannya singleton.

\textbf{25.} (i) Konstruksinya memakai: kestabilan \omterm{def:b1:topology:closed}{himpunan tertutup}
di bawah irisan sebarang (yaitu keberadaan $C$ sebagai \omterm{def:b1:topology:closed}{himpunan
tertutup}), teorema kekompakan \cref{thm:b1:topology:compact}
(pertanyaan 2, 22, 23), dan pencirian ketertutupan maupun
kelekatan lewat barisan (yaitu argumen kompak bersarang pada
\cref{ex:b1:topology:nested} dan pertanyaan 23). (ii) Keempat
pasangannya: panjang nol namun takterbilang (pertanyaan 3, 9); \omterm{def:b1:topology:closed}{tertutup}
namun berinterior kosong (pertanyaan 4); sempurna --- tanpa titik
yang terasing --- namun terputus seluruhnya (pertanyaan 11, 13);
serta terabaikan namun dengan $C + C = \intcc{0}{2}$ (pertanyaan 16). (iii)
Adapun surjeksi $h$ pada pertanyaan 20, bila dibuat kontinu dan
tidak turun, menjadi tangga setan pada teori
fungsi yang kontinu; sedangkan pada teori ukuran jilid Tahun ke-3,
$C$ menjadi saksi baku bahwa ``terabaikan'' tak berarti
``terbilang'', dengan sepupunya yang gemuk (pertanyaan 22) memisahkan
``tak \omterm{def:b1:topology:dense}{padat} di mana pun'' dari ``terabaikan''.

\end{solution}
